\documentclass[10pt,a4paper]{amsart}
\usepackage{amsmath,amssymb,amsthm,graphicx}
% Retained minimal preamble for the Stacks Project chapter sites-cohomology.tex (Cohomology on Sites).
% Verbatim excerpt of the chapter preamble: Xy-pic setup (native preamble lines 14-18) and the native
% theorem declarations with the mathematical macros (native preamble lines 154-217). Copied mechanically;
% nothing added, removed or redefined. The wrapper class and the amsmath/amssymb/amsthm/graphicx base
% packages are supplied by the assembler, so the class/externaldocument part of the native preamble is omitted.
% For commutative diagrams we use Xy-pic
\usepackage[all]{xy}

% We use 2cell for 2-commutative diagrams.
\xyoption{2cell}
% Theorem environments.
%
\theoremstyle{plain}
\newtheorem{theorem}[subsection]{Theorem}
\newtheorem{proposition}[subsection]{Proposition}
\newtheorem{lemma}[subsection]{Lemma}

\theoremstyle{definition}
\newtheorem{definition}[subsection]{Definition}
\newtheorem{example}[subsection]{Example}
\newtheorem{exercise}[subsection]{Exercise}
\newtheorem{situation}[subsection]{Situation}

\theoremstyle{remark}
\newtheorem{remark}[subsection]{Remark}
\newtheorem{remarks}[subsection]{Remarks}

\numberwithin{equation}{subsection}

% Macros
%
\def\lim{\mathop{\mathrm{lim}}\nolimits}
\def\colim{\mathop{\mathrm{colim}}\nolimits}
\def\Spec{\mathop{\mathrm{Spec}}}
\def\Hom{\mathop{\mathrm{Hom}}\nolimits}
\def\Ext{\mathop{\mathrm{Ext}}\nolimits}
\def\SheafHom{\mathop{\mathcal{H}\!\mathit{om}}\nolimits}
\def\SheafExt{\mathop{\mathcal{E}\!\mathit{xt}}\nolimits}
\def\Sch{\mathit{Sch}}
\def\Mor{\mathop{\mathrm{Mor}}\nolimits}
\def\Ob{\mathop{\mathrm{Ob}}\nolimits}
\def\Sh{\mathop{\mathit{Sh}}\nolimits}
\def\NL{\mathop{N\!L}\nolimits}
\def\CH{\mathop{\mathrm{CH}}\nolimits}
\def\proetale{{pro\text{-}\acute{e}tale}}
\def\etale{{\acute{e}tale}}
\def\QCoh{\mathit{QCoh}}
\def\Ker{\mathop{\mathrm{Ker}}}
\def\Im{\mathop{\mathrm{Im}}}
\def\Coker{\mathop{\mathrm{Coker}}}
\def\Coim{\mathop{\mathrm{Coim}}}

% Boxtimes
%
\DeclareMathSymbol{\boxtimes}{\mathbin}{AMSa}{"02}

%
% Macros for moduli stacks/spaces
%
\def\QCohstack{\mathcal{QC}\!\mathit{oh}}
\def\Cohstack{\mathcal{C}\!\mathit{oh}}
\def\Spacesstack{\mathcal{S}\!\mathit{paces}}
\def\Quotfunctor{\mathrm{Quot}}
\def\Hilbfunctor{\mathrm{Hilb}}
\def\Curvesstack{\mathcal{C}\!\mathit{urves}}
\def\Polarizedstack{\mathcal{P}\!\mathit{olarized}}
\def\Complexesstack{\mathcal{C}\!\mathit{omplexes}}
% \Pic is the operator that assigns to X its picard group, usage \Pic(X)
% \Picardstack_{X/B} denotes the Picard stack of X over B
% \Picardfunctor_{X/B} denotes the Picard functor of X over B
\def\Pic{\mathop{\mathrm{Pic}}\nolimits}
\def\Picardstack{\mathcal{P}\!\mathit{ic}}
\def\Picardfunctor{\mathrm{Pic}}
\def\Deformationcategory{\mathcal{D}\!\mathit{ef}}
\usepackage{cleveref}
\crefname{originallemma}{original lemma}{original lemmas}
\Crefname{originallemma}{Original lemma}{Original lemmas}
\makeatletter
\newlabel{lemma-higher-direct-images}{{072W}{source}{}{sourcepointer-1}{}}
\newlabel{lemma-higher-direct-images@cref}{{[originallemma][0][]072W}{[1][0][]source}{}{}{}}
\makeatother
\crefname{originallemma}{original lemma}{original lemmas}
\Crefname{originallemma}{Original lemma}{Original lemmas}
\makeatletter
\newlabel{sites-section-sheafification}{{00W1}{source}{}{sourcepointer-2}{}}
\newlabel{sites-section-sheafification@cref}{{[originallemma][0][]00W1}{[1][0][]source}{}{}{}}
\makeatother
\crefname{originallemma}{original lemma}{original lemmas}
\Crefname{originallemma}{Original lemma}{Original lemmas}
\makeatletter
\newlabel{categories-lemma-adjoint-exact}{{0038}{source}{}{sourcepointer-3}{}}
\newlabel{categories-lemma-adjoint-exact@cref}{{[originallemma][0][]0038}{[1][0][]source}{}{}{}}
\makeatother
\crefname{originallemma}{original lemma}{original lemmas}
\Crefname{originallemma}{Original lemma}{Original lemmas}
\makeatletter
\newlabel{sites-lemma-isomorphism-sheafifications}{{0H6Y}{source}{}{sourcepointer-4}{}}
\newlabel{sites-lemma-isomorphism-sheafifications@cref}{{[originallemma][0][]0H6Y}{[1][0][]source}{}{}{}}
\makeatother
\title{Higher direct images commute with filtered colimits over a covering family of objects}
\author{The Stacks Project Authors}
\date{stacks-project:sites-cohomology.tex:local-snapshot-20261011; GNU FDL 1.2 or later}
\begin{document}
\maketitle
\noindent\textit{Editorial note.} All delivered mathematical text is the author's own, copied from the pinned Stacks Project chapter file sites-cohomology.tex; no statement or proof step was rewritten, reordered or removed.\par\medskip
\noindent\textit{Editorial note.} The proof is the author's complete argument; the reduction and the application of the hypothesis are both written out.\par\medskip
\noindent\textit{Editorial note.} Every labelled reference occurring inside the counted spans is carried in assembly.external\_references as a non-load-bearing navigation pointer (same-chapter ones at their definition site, cross-chapter ones at the invocation inside the counted span), so no load-bearing mathematics is replaced by a stub. No bibliography citation occurs inside the counted spans.\par\medskip
\noindent\textit{Editorial note.} The checked chapter list for same-name collisions is the set of stacks chapters staged in this repository: chow, cohomology, divisors, etale-cohomology, flat, more-algebra, more-morphisms, morphisms, spaces-more-morphisms, varieties together with this card's chapter sites-cohomology.\par\medskip
\noindent\textit{Editorial note.} The preamble is the chapter's own retained minimal preamble (sources/stacks-sites-cohomology/retained-preamble.tex), a verbatim excerpt of the native preamble: the Xy-pic setup and the native theorem declarations with the mathematical macros.\par\medskip
\noindent\textit{Allowed standard prerequisites (editorial).}\begin{itemize}
\item Morphism of sites and higher direct images of abelian sheaves. Same chapter, the functors f\_*, R\textasciicircum{}pf\_* attached to a morphism of sites and their description by cohomology on the objects of the target site. Supplies the functors appearing in the statement.
\item Filtered colimits of abelian sheaves and the commutation of cohomology with such colimits on the objects of B. Same chapter, the criterion for cohomology to commute with filtered colimits, which is exactly the hypothesis on the set B in the statement. Supplies the reduction from the sheaf statement to the pointwise one.
\end{itemize}\medskip
\ifdefined\hypertarget\hypertarget{sourcepointer-1}{}\fi
\noindent\textit{Source reference.} Original lemma 072W: sources/stacks-sites-cohomology/pinned-original-source.tex, source lines 747--747. Navigation pointer invoked at proof line 2837; the counted span contains the pointer only and the cited result is not reproduced inside it, so the pointer is registered as a non-load-bearing navigation stub in the established form for Stacks material. No bibliography citation is involved.\par\medskip
\ifdefined\hypertarget\hypertarget{sourcepointer-2}{}\fi
\noindent\textit{Source reference.} Original lemma 00W1: sources/stacks-sites-cohomology/pinned-original-source.tex, source lines 2842--2842. Navigation pointer invoked at proof line 2842; the counted span contains the pointer only and the cited result is not reproduced inside it, so the pointer is registered as a non-load-bearing navigation stub in the established form for Stacks material. No bibliography citation is involved.\par\medskip
\ifdefined\hypertarget\hypertarget{sourcepointer-3}{}\fi
\noindent\textit{Source reference.} Original lemma 0038: sources/stacks-sites-cohomology/pinned-original-source.tex, source lines 2843--2843. Navigation pointer invoked at proof line 2843; the counted span contains the pointer only and the cited result is not reproduced inside it, so the pointer is registered as a non-load-bearing navigation stub in the established form for Stacks material. No bibliography citation is involved.\par\medskip
\ifdefined\hypertarget\hypertarget{sourcepointer-4}{}\fi
\noindent\textit{Source reference.} Original lemma 0H6Y: sources/stacks-sites-cohomology/pinned-original-source.tex, source lines 2850--2850. Navigation pointer invoked at proof line 2850; the counted span contains the pointer only and the cited result is not reproduced inside it, so the pointer is registered as a non-load-bearing navigation stub in the established form for Stacks material. No bibliography citation is involved.\par\medskip
\begin{lemma}
\label{lemma-higher-direct-image-colimit}
Let $f : \mathcal{C} \to \mathcal{D}$ be a morphism of sites
corresponding to the continuous functor $u : \mathcal{D} \to \mathcal{C}$.
Let $(\mathcal{F}_i, \varphi_{ii'})$ be a system of
abelian sheaves on $\mathcal{C}$. Set $\mathcal{F} = \colim \mathcal{F}_i$.
Let $p \geq 0$ be an integer. Denote $\mathcal{B}$ the set of
$V \in \Ob(\mathcal{D})$ such that
$H^p(u(V), \mathcal{F}) = \colim H^p(u(V), \mathcal{F}_i)$.
If every object of $\mathcal{D}$ has a covering by elements of $\mathcal{B}$,
then $R^pf_*\mathcal{F} = \colim R^pf_*\mathcal{F}_i$.
\end{lemma}

\begin{proof}
Recall that $R^pf_*\mathcal{F}$ is the sheafification of the presheaf
$\mathcal{G}$ sending $V$ to $H^p(u(V), \mathcal{F})$, see
Lemma \ref{lemma-higher-direct-images}. Similarly, 
$R^pf_*\mathcal{F}_i$ is the sheafification of the presheaf
$\mathcal{G}_i$ sending $V$ to $H^p(u(V), \mathcal{F}_i)$.
Recall that sheafification is the left adjoint to the inclusion
from sheaves to presheaves, see Sites, Section
\ref{sites-section-sheafification}. Hence sheafification commutes
with colimits, see Categories, Lemma \ref{categories-lemma-adjoint-exact}.
Hence it suffices to show that the map of presheaves (with colimit
in the category of presheaves)
$$
\colim \mathcal{G}_i \longrightarrow \mathcal{G}
$$
induces an isomorphism on sheafifications. This follows from
Sites, Lemma \ref{sites-lemma-isomorphism-sheafifications} and
our assumption on $\mathcal{B}$.
\end{proof}

\end{document}
